\documentclass[11pt]{article}
\usepackage[margin=1in]{geometry}
\usepackage{listings,xcolor,hyperref,booktabs}
\lstset{basicstyle=\ttfamily\small,breaklines=true,frame=single,columns=fullflexible,
  keywordstyle=\color{blue},commentstyle=\color{gray},showstringspaces=false,
  morekeywords={fn,let,mut,match,struct,impl,const,static,use,pub,loop,while,for,in,if,else,move,macro_rules}}
\title{A Tiny Media Server in Rust\\\large Notes for a C++ Programmer Who Has Never Seen Rust}
\date{}
\emergencystretch=3em
\begin{document}
\maketitle
\tableofcontents

\section{What we built}
A program that listens on a TCP port and answers every connection with one audio (or video) file baked into the executable. Open \texttt{http://localhost:3005} in a browser and it plays.

\begin{itemize}
\item \texttt{src/main.rs} -- the whole server (about 80 lines).
\item \texttt{media/3005.bbd} -- the file that gets embedded. Replace it and rebuild.
\item \texttt{Cargo.toml} -- project and optimisation settings.
\item \texttt{.cargo/config.toml} -- makes the default build target static Linux (musl).
\end{itemize}

Build and run:
\begin{lstlisting}[language=bash]
cargo build --release
./target/x86_64-unknown-linux-musl/release/mediaserver
\end{lstlisting}

\section{Rust from the ground up (for C++ people)}
\subsection{Tooling}
\texttt{cargo} is CMake + a package manager + a test runner in one. \texttt{rustc} is the compiler (like \texttt{g++}); you almost never call it directly. \texttt{cargo build --release} is an optimised build; plain \texttt{cargo build} is a debug build.

\subsection{Basic syntax}
\begin{lstlisting}
fn add(a: i32, b: i32) -> i32 {   // types come after names
    a + b                          // no semicolon = this is the return value
}
let x = 5;          // immutable by default (like const auto)
let mut y = 5;      // mutable
const PORT: u16 = 3005;   // compile-time constant, type required
\end{lstlisting}
Integer types are explicit: \texttt{u8, u16, u32, u64, usize} (unsigned; \texttt{usize} is pointer-sized, like \texttt{size\_t}) and \texttt{i8..i64, isize} (signed). Almost everything is an expression: \texttt{if}, \texttt{match} and blocks all produce values.

\subsection{Ownership and borrowing}
This is the big difference from C++. Every value has exactly one owner; when the owner goes out of scope the value is dropped (RAII, always). Assigning or passing a value \emph{moves} it -- the old name becomes unusable (like \texttt{std::move}, but enforced by the compiler). Instead of moving you can \emph{borrow}:

\begin{center}\begin{tabular}{ll}\toprule
Rust & Roughly C++ \\\midrule
\texttt{\&T} & \texttt{const T\&} (many allowed at once) \\
\texttt{\&mut T} & \texttt{T\&} (exactly one, and no \texttt{\&T} at the same time) \\
\texttt{Vec<u8>} & \texttt{std::vector<uint8\_t>} \\
\texttt{\&[u8]} & \texttt{std::span<const uint8\_t>} (pointer + length, a ``slice'') \\
\texttt{String} / \texttt{\&str} & \texttt{std::string} / \texttt{std::string\_view} \\
\texttt{Box<T>} & \texttt{std::unique\_ptr<T>} \\
\texttt{Option<T>} & \texttt{std::optional<T>} \\
\texttt{Result<T,E>} & value-or-error (no exceptions in Rust) \\
\bottomrule\end{tabular}\end{center}

The compiler proves at build time that no reference outlives what it points to and that you never have a data race. That removes dangling pointers and use-after-free in safe code.

\subsection{Lifetimes and \texttt{'static}}
A reference type can carry a lifetime: \texttt{\&'a [u8]}. \texttt{\&'static} means ``valid for the whole program''. String literals and \texttt{include\_bytes!} data are \texttt{'static} because they live in the executable's read-only data.

\subsection{Errors: \texttt{Result}, \texttt{?}, \texttt{Option}}
\begin{lstlisting}
let n = s.read(&mut buf)?;   // if Err, return it from this function immediately
let v = x.expect("msg");     // if Err/None, crash with msg (fine at startup)
let _ = might_fail();        // deliberately ignore a result
\end{lstlisting}
Ignoring a \texttt{Result} without \texttt{let \_ =} produces a compiler warning -- Rust makes you say that you meant it.

\subsection{Structs, modules, macros}
\texttt{struct Resp \{ ... \}} is a plain data type (no inheritance). Names ending in \texttt{!} (\texttt{include\_bytes!}, \texttt{format!}, \texttt{macro\_rules!}) are macros, expanded at compile time -- hygienic, and far saner than the C preprocessor. \texttt{use std::io::Read;} is like \texttt{using}/\texttt{\#include}: it brings a name (here a \emph{trait}, similar to an interface) into scope so its methods can be called.

\subsection{Closures and threads}
\texttt{|a, b| a + b} is a lambda. \texttt{move || ...} captures by value, i.e.\ the closure takes ownership of what it uses. \texttt{std::thread::spawn(closure)} runs it on a new OS thread. The compiler only allows data to cross into a thread if it is safe to do so -- here, either owned by the closure or \texttt{'static}.

\section{Walking through \texttt{main.rs}}
\subsection{Constants and embedding the file}
The port is hardcoded to 3005. The line \texttt{const \_: () = assert!(PORT == 0x0BBD, ...)} is evaluated by the compiler: \texttt{0x0BBD} is 3005 in hexadecimal ($11\cdot256+11\cdot16+13$), so editing \texttt{PORT} to any other value is a \emph{build error}, not a runtime one. (\texttt{const \_} is an unnamed constant whose only job is to run the check, similar to a C++ \texttt{static\_assert}.)

\begin{lstlisting}
const PORT: u16 = 3005;
macro_rules! media_path { () => { "../media/3005.bbd" }; }
const PATH: &str = media_path!();
static MEDIA: &[u8] = include_bytes!(media_path!());
\end{lstlisting}
\texttt{include\_bytes!} reads the file at \emph{compile time} and puts its bytes in the executable (think \texttt{xxd -i}, built in). Its argument must be a string literal, not a \texttt{const}, so a tiny macro supplies the path in one place; the same text is reused as \texttt{PATH} to choose the MIME type. The path is relative to the source file (\texttt{src/}), hence \texttt{../media}.

\subsection{Building one buffer at startup}
We format the header once, append the body, then \texttt{Box::leak} the buffer. \texttt{Box::leak} deliberately never frees the memory and hands back a \texttt{\&'static mut}; we use it as an immutable \texttt{\&'static [u8]}. Leaking is fine because the data should live as long as the process, and being \texttt{'static} lets every thread share it with no locks and no reference counting.

The header is:
\begin{lstlisting}
HTTP/1.1 200 OK\r\n
Content-Type: audio/mpeg\r\n
Content-Length: <n>\r\n
Accept-Ranges: bytes\r\n
Connection: close\r\n
\r\n
\end{lstlisting}
HTTP headers end with an empty line (\verb|\r\n\r\n|). \texttt{Content-Length} tells the browser when the body ends; \texttt{Connection: close} says we won't reuse the socket. \texttt{Accept-Ranges} advertises seeking support (an addition to the minimum list; Safari likes it).

\subsection{The accept loop}
\begin{lstlisting}
let listener = TcpListener::bind(("0.0.0.0", PORT)).expect("bind");
for stream in listener.incoming().flatten() {
    std::thread::spawn(move || { let _ = handle(stream, resp); });
}
\end{lstlisting}
\texttt{TcpListener} wraps \texttt{socket/bind/listen/accept}. \texttt{incoming()} yields a \texttt{Result<TcpStream>} for each connection; \texttt{.flatten()} drops the failures. Each accepted stream is \emph{moved} into a fresh thread. When \texttt{handle} returns, the stream is dropped, which closes the socket (RAII).

\subsection{Reading the request}
We loop on \texttt{read} until the bytes seen so far contain \verb|\r\n\r\n| (capped at 64 KiB so a rude client can't make us buffer forever). We never parse the method or path -- every URL gets the same file.

\subsection{Range requests}
When you seek, or when Safari/iOS opens media, the browser sends \texttt{Range: bytes=100-199}. We answer \texttt{206 Partial Content} with \texttt{Content-Range: bytes 100-199/total} and only those bytes. \texttt{parse\_range} accepts a single \texttt{bytes=a-b} or \texttt{bytes=a-}; anything else (multiple ranges, suffix form \texttt{-500}, garbage, out-of-bounds start) returns \texttt{None} and we fall back to the full 200. The end is clamped to the last byte. A 206 header depends on the request so it is formatted per request (the 200 path remains one pre-built \texttt{write\_all}); the body slice is zero-copy: \texttt{\&body[a..=b]}.

\subsection{Why drain before closing?}
If a program closes a TCP socket while unread data is still in its receive buffer, Linux sends a \emph{RST} (reset) instead of a normal \emph{FIN}. A RST can make the client discard response bytes that were still in flight, truncating the file. So after writing we call \texttt{shutdown(Write)} (sends FIN -- ``I'm done sending'') and then read until the client closes its side, and only then drop the socket.

\section{Cargo.toml and the build}
\begin{center}\begin{tabular}{lp{0.6\textwidth}}\toprule
Setting & Effect \\\midrule
\texttt{opt-level = 3} & like \texttt{-O3} \\
\texttt{lto = "fat"} & whole-program link-time optimisation \\
\texttt{codegen-units = 1} & one unit: slower compile, better optimisation \\
\texttt{panic = "abort"} & a panic kills the process (no unwinding code); a panic in one connection thread therefore takes down the whole server \\
\texttt{strip = true} & remove symbols; smaller binary \\
\bottomrule\end{tabular}\end{center}
The target \texttt{x86\_64-unknown-linux-musl} links against musl libc statically, so the binary has no dependencies and runs on any x86-64 Linux. One-time setup: \texttt{rustup target add x86\_64-unknown-linux-musl}. No \texttt{libc} crate or \texttt{sendfile} was needed; with the file in memory, \texttt{write\_all} is already a straight copy to the kernel.

\section{Verifying it}
\begin{lstlisting}[language=bash]
curl -I localhost:3005                  # headers
curl -r 0-99 -o part.bin localhost:3005 # 206, exactly 100 bytes
curl -o all.bin localhost:3005
cmp all.bin media/3005.bbd            # no output = identical
\end{lstlisting}
\texttt{curl -I} sends a HEAD request; because we do not parse methods the server also sends the body, which curl ignores.

\section{Security review}
Threat model: anyone who can reach port 3005 sends arbitrary bytes. Goal: they can't run code, escalate, or kill the server.
\begin{itemize}
\item \textbf{No memory-unsafe code.} \texttt{\#![forbid(unsafe\_code)]} makes the compiler reject any \texttt{unsafe} block, and there are no dependencies, so buffer overflows and use-after-free are not possible in this program.
\item \textbf{Request is never interpreted.} Only the \texttt{Range} header is read, and its numbers are parsed with overflow-checked \texttt{parse} and clamped to the file length. No input is ever echoed into a response (no header injection), used as a file path (no traversal), or passed to a shell. The audio is embedded, so the server never opens files at runtime.
\item \textbf{Panic = total outage.} With \texttt{panic = "abort"} any panic kills the whole server. Two panic sources were removed: slicing a \texttt{\&str} at a non-character boundary (a \texttt{Range:} line containing multibyte UTF-8 would have crashed the process, so parsing is now done on raw bytes) and \texttt{thread::spawn} failing (now \texttt{Builder::spawn}, whose error is handled).
\item \textbf{Resource limits.} Max 256 simultaneous connections (extras are closed immediately), 16\,KiB request cap, 10\,s read/write timeouts, bounded post-response drain, 64\,KiB thread stacks, and a header scan that is linear (not quadratic) in slow-drip input.
\item \textbf{Privilege.} The server refuses to start as root (it checks the effective uid in \texttt{/proc/self/ status}). Port 3005 is above 1024, so root is never needed. Run it as an unprivileged user, ideally under systemd with \texttt{DynamicUser=yes}, \texttt{NoNewPrivileges=yes}, \texttt{ProtectSystem=strict}, \texttt{ProtectHome=yes}, \texttt{PrivateTmp=yes}. The static binary needs no files at all.
\item \textbf{Audio-specific.} The server treats the file as opaque bytes and never decodes it, so a malformed MP3 cannot exploit the server. It could only target the \emph{listener's} decoder, so embed only audio you trust. Responses carry \texttt{X-Content-Type-Options: nosniff} (the browser must treat the bytes as audio, never sniff HTML/JS out of them) and \texttt{Content-Security-Policy: sandbox}.
\item \textbf{Not covered.} No TLS (traffic is plaintext, visible and modifiable on the network), and bandwidth floods must be handled upstream. No software can honestly be called ``provably'' secure; this is a small, dependency-free design reviewed and tested against the cases above.
\end{itemize}

\section{Limitations and exercises}
\begin{itemize}
\item A client that reads very slowly can hold one of the 256 connection slots for a long time (each individual read/write has a 10\,s timeout, but a whole transfer does not). Put a reverse proxy or firewall rate limit in front for hostile networks.
\item No HEAD, 404, 416 handling -- by design.
\item Exercise: support the suffix form \texttt{bytes=-500}.
\item Exercise: switch to \texttt{.mp4} (change the macro, drop in a file) and check the \texttt{Content-Type}.
\end{itemize}
\end{document}
